\documentclass[titlepage,a4paper]{jsarticle}
\usepackage{../../../sty/import}% 各種パッケージインポート
\usepackage{../../../sty/title}% タイトルページの変更
\usepackage{../../../sty/answergrid}
%% タイトルページの変数
% レポートタイトル
\title{心理学特論第13回目出席票}
% 提出日
\expdate{\today}
% 科目名
\subject{心理学特論}
% 分野
\class{情報経営システム工学分野}
% 学年
\grade{M1}
% 学籍番号
\mynumber{24336488}
% 記述者
\author{本間三暉}

\begin{document}
% titleページ作成
\maketitle
\section*{keyword\#1}
偽記憶
\section*{keyword\#2}
既視感
\section*{心理学特論第13回の主なテーマである「記憶」にまつわる悩みはいろいろありそうです。
「学習した内容が試験で思い出せない…」，「以前会った人の顔や名前が思い出せない…」など。
今回は「記憶」に関する日頃の心がけや経験に焦点を当てて，思い出しながら書いてみましょう。}
\section{きちんと記憶するために日頃から心がけていることがあれば，簡潔に書いて下さい。}
就職活動では，面接日程や提出期限，面接で改善したい点などを忘れないように，ToDoリストやメモアプリに記録し定期的に見返すようにしている．

\section{「忘れた！」，すなわち忘却にまつわる失敗のエピソードがあれば，簡潔に書いて下さい。}
就職活動で，Webテストの締切を当日の23時59分だと思い込んでいたため，実際の13時締切に間に合わず提出できなかった．
この経験から，締切日時は必ずメモやカレンダーで確認するようにしている．

\newpage
\section{心理学特論課題サイコロワーク10解答記入欄}

\AnswerGridFilled{
  {イ}
  {ウ}
  {ア}
  {エ}
  {エ}
  {二重貯蔵モデル}
  {短期記憶}
  {長期記憶}
  {宣言的記憶}
  {手続き的記憶}
  {長期記憶}
  {意味記憶}
  {エピソード記憶}
  {未来記憶}
  {既視感}
  {エビングハウス}
  {ヒグビー}
  {組織化}
  {ロフタス}
  {35倍}
}
\end{document}